\par\Needspace{4\baselineskip}
\begin{center}
\SourceLeaf{042}\pdfbookmark[2]{24.}{article-024}24.
\end{center}
Will man die Rechnung mit hyperbolischen Logarithmen durchführen, so verwendet man zweckmäßigerweise die Hilfsgröße $F$, die sich aus Gleichung III bestimmen lässt. Daraus erhält man $N$ nach Formel XI. Mit Formel VIII berechnet man den Halbparameter aus dem Radiusvektor oder umgekehrt den Radiusvektor aus dem Halbparameter. Den zweiten Teil von $N$ kann man, wenn man möchte, auf zwei Arten bestimmen, nämlich mit der Formel $\log\operatorname{hyp}\operatorname{tang}(45^0+\frac{1}{2}F)$ und mit $\log\operatorname{hyp}\cos\frac{1}{2}(v-\psi)-\log\operatorname{hyp}\cos\frac{1}{2}(v+\psi)$. Offenbar ist die Größe $N$ hier, wo $\lambda=1$ gilt, im Verhältnis $1:\lambda$ größer als bei Verwendung \textsc{Briggs}scher Logarithmen. Mit diesem Verfahren ergibt sich für unser Beispiel die folgende Rechnung.
\begin{center}
\begin{tabular}{@{}ll@{}}
\makebox[18em][l]{$\log\operatorname{tang}\frac{1}{2}\psi$\hfill .\,.\,.\,.\,.\,.\,.\,.\,.} & $9{,}531\,8179$\\[.3em]
\makebox[18em][l]{$\log\operatorname{tang}\frac{1}{2}v$\hfill .\,.\,.\,.\,.\,.\,.\,.\,.} & $9{,}220\,1009$\\[.3em]
\hline
\makebox[18em][l]{$\log\operatorname{tang}\frac{1}{2}F$\hfill .\,.\,.\,.\,.\,.\,.\,.\,.} & $8{,}751\,9188$\\[.3em]
\hline
\makebox[18em][l]{$\log e$\hfill .\,.\,.\,.\,.\,.\,.\,.\,.\,.\,.\,.\,.\,.} & $0{,}101\,0188$\\[.3em]
\makebox[18em][l]{$\log\operatorname{tang}F$\hfill .\,.\,.\,.\,.\,.\,.\,.\,.\,.} & $9{,}054\,3366$\\[.3em]
\hline
 & $9{,}155\,3554$\\[.3em]
$e\operatorname{tang}F=$ & $0{,}143\,0063\,8$\\[.3em]
$\log\operatorname{hyp}\operatorname{tang}(45^0+\frac{1}{2}F)=$ & $0{,}113\,0866\,6$\\[.3em]
\hline
$N=$ & $0{,}029\,9197\,2$
\end{tabular}
\end{center}
\begin{center}
$\frac{1}{2}F=3^0\,13^{\prime}\,58^{\prime\prime}{,}12$
\end{center}
\begin{center}
\begin{tabular}{@{}ll@{}}
K. $\log\operatorname{hyp}\cos\frac{1}{2}(v-\psi)=$ & $0{,}013\,4226\,6$\\[.3em]
K. $\log\operatorname{hyp}\cos\frac{1}{2}(v+\psi)=$ & $0{,}126\,5093\,0$\\[.3em]
\hline
Diff. $=$ & $0{,}113\,0866\,4$\\[.7em]
\hline
\makebox[18em][l]{$\log N$\hfill .\,.\,.\,.\,.\,.\,.\,.\,.\,.\,.\,.\,.} & $8{,}475\,9575$
\end{tabular}
\end{center}
\begin{center}
\footnotesize
\begin{tabular}{@{}ll@{}}
$\left.\begin{array}{ll}
\log k\quad .\,.\,.\,.\,.\,.\,.\,.\,.\,.\,.\,.\,.\,. &8{,}235\,5814\\[.3em]
\frac{3}{2}\log b\quad .\,.\,.\,.\,.\,.\,.\,.\,.\,.\,.\,.\,. &0{,}903\,0900
\end{array}\right\}$ & \begin{tabular}{@{}ll@{}}
Diff. .\,.\,.\,.\,.\,.\,.\,.\,.\,.\,.\,.\,. & $7{,}332\,4914$\\[.3em]
\hline
$\log t$ .\,.\,.\,.\,.\,.\,.\,.\,.\,.\,.\,.\,.\,. & $1{,}143\,4661$\\[.3em]
$t=$ & $13{,}91445$
\end{tabular}
\end{tabular}
\end{center}
